\chapter{El trabajo que se fue}
\label{cap:trabajo}

El capítulo anterior enunció lo que los otros documentan sobre el agua. Este hace lo mismo con
el trabajo, y encuentra que el libro también tiene una economía escrita sin haberla enunciado:
tabaco en tres capítulos, canteras y áridos en seis, turismo y alojamientos en otros seis,
aserradero, cooperativas, productores organizados.

Como el anterior, \textbf{sobre todo pone en orden}, aunque suma el procesamiento censal de 2022 y tres sociedades recientes.

\section{La frase}

\textbf{Este departamento dejó de producir alimentos, empezó a producir insumos para otros, y
el Estado alentó por decreto uno de esos movimientos.}

Los tres movimientos están documentados por separado y ninguno se leyó junto a los demás.

\section{Primero: la desaparición}

El dato es brutal y está en el capítulo \ref{cap:tierra}. Entre 2002 y 2018 el departamento
pasó de \textbf{308 explotaciones agropecuarias a 131}, y de \textbf{98.774,9 hectáreas a
11.417,5}: una caída del \textbf{88,4 por ciento} de la superficie agropecuaria en dieciséis
años.

El capítulo verifica que no es un artefacto censal. En el mismo período la superficie
agropecuaria de la provincia \textbf{creció} un 2,8 por ciento. En La Caldera no hubo
concentración: hubo desaparición.

Y hubo, antes, producción reconocida. Un decreto de enero de 1980 declara zona de emergencia
agropecuaria por daños de granizo e incluye expresamente \textbf{el área tabacalera} de La
Caldera. En 1934 la provincia había incorporado al departamento a la «región económica del
olivo», con exención impositiva por diez años para quien plantara olivos. Hubo un
\textbf{aserradero} en el pueblo.

\textbf{Nada de eso figura hoy en el censo agropecuario.}

\begin{cautela}
El capítulo \ref{cap:tierra} es explícito en algo que conviene repetir acá, porque la
tentación de explicarlo todo con una causa es fuerte.

\textbf{La hipótesis de que el colapso del tabaco explique la pérdida de superficie quedó
refutada para 2002--2018} por el propio dato censal: los cultivos industriales crecieron, y el tabaco no tenía peso suficiente en este departamento
para arrastrar 87.357 hectáreas. Lo que el decreto de 1980 prueba es que había área tabacalera
reconocida, no que su desaparición explique la caída.

Qué la explica, este libro lo sabe sólo en parte: el censo de 2008 indica que buena parte de esa superficie era monte y pastizal que dejó de declararse (capítulo \ref{cap:tierra}).
\end{cautela}

\section{Segundo: lo que ocupó el lugar}

Dos actividades crecieron mientras la agropecuaria desaparecía, y las dos tienen una
característica común que conviene nombrar.

\textbf{La extracción de áridos.} El capítulo \ref{cap:resistencias} documenta veinte
concesiones vigentes sobre los cauces del corredor, seis de ellas en manos de una sola
cooperativa. El material no se consume en el departamento: alimenta la obra de la capital y de
la región.

\textbf{La vivienda de fin de semana.} El capítulo \ref{cap:poblacion} establece que el
departamento tiene una proporción de viviendas sin residentes permanentes que \textbf{duplica
la provincial} ---8,55 por ciento en la provincia--- y que contiene una fracción censal con el
\textbf{43,17 por ciento} de sus viviendas sin hogar residente, entre las cinco más altas de Salta.

\textbf{Lo que las dos tienen en común es a quién pertenecen.} La piedra sale hacia Salta; las
casas vacías son de gente de Salta. El capítulo \ref{cap:aguabaja} mostró que el agua hace el
mismo recorrido. \textbf{Son tres flujos en la misma dirección, y ninguno de los tres fue
decidido como política.}

\subsection*{Y el empleo no cayó: el capítulo \ref{cap:poblacion} ya lo había establecido}

Sería esperable que la desaparición de la tierra agropecuaria hubiera dejado desocupación rural.
No la dejó, y este libro lo sabía antes de escribir este capítulo.

El capítulo \ref{cap:poblacion} procesó el Censo 2010 y encontró que La Caldera tiene
\textbf{la tasa de ocupación más alta de los veintitrés departamentos de Salta} ---64,37 por
ciento de ocupados sobre población de catorce años y más, contra 54,42 provincial--- y
\textbf{la tasa de inactividad más baja}. Y concluyó, juntando eso con los hogares chicos y la
vivienda vacía, que el perfil es el de una \textbf{ciudad dormitorio}.

\textbf{Este capítulo agrega tres cosas a eso.}

\begin{ficha}{Censo 2022 --- mismo procesamiento, doce años después}
\textbf{La Caldera: actividad 68,97 --- empleo 64,71 --- desocupación 6,18.}

\textbf{Primero otra vez}, en las dos tasas, sobre los veintitrés departamentos. La sigue Cachi
con 65,80 y 63,28; la Capital queda tercera en actividad, con 64,39.

El promedio provincial se ubicó en 58,59 y 54,57, y la distancia se mantuvo: en empleo, el departamento sigue \textbf{unos diez puntos por encima}, como en 2010.

Su desocupación creció 1,39 puntos entre los dos censos. La de la Capital creció 0,94 en el
mismo período: \textbf{subió en todas partes}.
\end{ficha}

\textbf{Lo primero, entonces: no es un dato de un censo, es una constante de doce años}, con la
superficie agropecuaria desplomándose en el medio.

\textbf{Lo segundo: adentro del departamento las diferencias son menores.}

\begin{ficha}{Censo 2022, tasas por fracción censal}
\textbf{66077-01} ---la más poblada---: actividad 68,68, empleo 64,72, desocupación 5,77.

\textbf{66077-02}: actividad \textbf{70,33}, la más alta de las tres; empleo 65,49.

\textbf{66077-03} ---la de las viviendas vacías, con 43,17\,\% sin residentes permanentes---:
actividad 65,53, empleo \textbf{60,98}, la más baja de las tres.
\end{ficha}

Las tres están por encima del promedio provincial. Incluso la más baja supera en empleo a la
Capital, que registra 58,99. \textbf{Podría suponerse que la fracción de las casas de fin de
semana tuviera empleo bajo; tiene casi cuatro puntos menos que el promedio departamental y sigue siendo alta.}

\textbf{Y lo tercero, que es lo que importa para este capítulo.}

\begin{cautela}
\textbf{La caída del ochenta y ocho por ciento de la superficie agropecuaria no se tradujo en
desempleo, y eso vuelve la pregunta del capítulo \ref{cap:tierra} más urgente, no menos.}

Si el departamento perdió ochenta y siete mil hectáreas entre 2002 y 2018 y su tasa de empleo
no se movió, entonces quienes trabajaban esa tierra \textbf{o se fueron, o encontraron otra
cosa, o nunca fueron tantos como para notarse en la tasa}. Las tres son posibles y este libro
no sabe cuál.

\textbf{Y la lectura de ciudad dormitorio del capítulo \ref{cap:poblacion} sigue siendo una
inferencia, no una medición.} El censo pregunta dónde vive la persona, no dónde trabaja. Que la
base productiva local sea escasa hace razonable suponer que trabajan en la capital, pero
\textbf{suponer no es establecer}. La variable que lo resolvería ---lugar de trabajo--- no está
disponible en el procesamiento público y quedó pedida en el apéndice \ref{cap:pedidos}.

\textbf{Lo que sí queda retirado de este capítulo} es la presunción de que las actividades que
crecieron «dejan poco trabajo» en general: para el empleo del departamento no está probada, y la tasa apunta en contra. Para las dos actividades que crecieron, en cambio, el procesamiento por rama que sigue la confirma.
\end{cautela}

\subsection*{Y de qué viven: la rama de actividad}

El mismo censo permite, por fin, contestar la pregunta que este capítulo venía rodeando.

\begin{ficha}{Censo 2022 --- rama de actividad de la población ocupada, La Caldera}
Sobre \textbf{6.253 ocupados}, y en orden:

\textbf{Comercio} 733 ---11,72\,\%---. \textbf{Construcción} 678 ---10,84\,\%---.
\textbf{Enseñanza} 505 ---8,08\,\%---. \textbf{Personal doméstico en hogares} 464
---7,42\,\%---. \textbf{Administración pública, defensa y seguridad social} 347 ---5,55\,\%---.
\textbf{Industria manufacturera} 309 ---4,94\,\%---. \textbf{Salud} 297 ---4,75\,\%---.
\textbf{Servicios profesionales, científicos y técnicos} 251 ---4,01\,\%---.
\textbf{Alojamiento y servicios de comida} \textbf{218} ---3,49\,\%---.
\textbf{Agricultura, ganadería, caza, silvicultura y pesca} \textbf{193} ---3,09\,\%---.
Transporte 138. \textbf{Explotación de minas y canteras} \textbf{51} ---0,82\,\%---.
Servicios inmobiliarios 25.

\textbf{Sin respuesta o con información insuficiente para codificar: 1.546 ---24,7\,\%---.}

Las ramas no enumeradas reúnen a las 498 personas restantes.
\end{ficha}

\textbf{Cuatro cosas se leen ahí, y dos obligan a corregir lo que este capítulo dijo antes.}

\textbf{Primera: el sector público y el educativo-sanitario son el mayor empleador.} Enseñanza,
administración pública y salud suman \textbf{1.149 personas, el 18,4 por ciento}. Cada una de
las tres está por encima del promedio provincial: enseñanza 8,08 contra 5,53; salud 4,75 contra
3,25; administración pública 5,55 contra 4,62.

\textbf{Segunda: el servicio doméstico emplea a 464 personas}, el 7,42 por ciento. Es el cuarto
rubro del departamento, \textbf{más del doble que la hotelería y la gastronomía juntas}, y está
exactamente en el promedio provincial.

\textbf{Tercera: la agricultura emplea a 193 personas y las canteras a 51.} El capítulo
\ref{cap:resistencias} documenta veinte concesiones de áridos vigentes sobre los cauces del
corredor; \textbf{ocupan a cincuenta y una personas}. Es la rama más chica de todas las
relevantes del departamento.

\textbf{Y cuarta: el turismo emplea 218 personas}, el 3,49 por ciento, \textbf{por debajo del
promedio provincial} ---3,84--- y por debajo de la Capital ---4,61---.

\begin{cautela}
\textbf{Esto obliga a corregir el tono de este capítulo, y en dos direcciones opuestas.}

\textbf{Hacia abajo, la lectura turística.} El capítulo describió una transformación en la que
el turismo y la vivienda de fin de semana ocupan el lugar que dejó la producción agropecuaria.
\textbf{En empleo, el turismo pesa menos que en el promedio de la provincia}, y menos que el
servicio doméstico, la construcción o la enseñanza.

\textbf{Hacia abajo también, la lectura minera.} Las veinte concesiones de áridos que este
libro documenta con expediente, canon e informe ambiental, y que extraen del cauce del río,
\textbf{ocupan a cincuenta y una personas}. Lo que el capítulo \ref{cap:aguabaja} describe como
un flujo que sale del departamento no es, en ningún sentido, una fuente de trabajo local.

\textbf{Y hacia arriba, la agropecuaria.} El capítulo \ref{cap:tierra} documenta que la
superficie agropecuaria cayó un ochenta y ocho por ciento. \textbf{Y todavía hay 193 personas
cuya actividad principal es agrícola o ganadera}, casi tantas como en el turismo. La actividad
no desapareció: desapareció la superficie declarada en el censo agropecuario. \textbf{Son dos
cosas distintas y este libro las había tratado como una sola.}
\end{cautela}

\begin{cautela}
\textbf{Un límite serio de estos números, y conviene ponerlo antes de que alguien lo use mal.}

\textbf{Uno de cada cuatro ocupados del departamento ---1.546 personas--- no tiene rama
codificada}: 825 sin respuesta y 721 con información insuficiente. Es una proporción parecida a
la provincial, de modo que no es una anomalía de La Caldera, pero sigue siendo un cuarto del
universo.

\textbf{Eso significa que cualquiera de las cifras anteriores podría ser mayor}, y que las ramas
más chicas ---canteras, inmobiliarias, transporte--- son las más sensibles a ese faltante.

\textbf{Y sigue sin resolverse dónde trabajan.} La rama dice qué hacen, no en qué jurisdicción.
Una parte de esas 733 personas en comercio y de esas 505 en enseñanza puede estar trabajando en
la capital. Es lo que la lectura de ciudad dormitorio del capítulo \ref{cap:poblacion} supone y
lo que este libro no puede establecer.
\end{cautela}

\begin{cautela}
\textbf{Una precisión sobre el procesamiento.} La consulta de 2022 se corrió sobre dos bases
distintas ---«viviendas particulares» y «pueblos originarios, afrodescendientes e identidad de
género»--- y \textbf{devolvió valores idénticos en las dos}. Eso indica que el indicador se
calcula sobre la población total con independencia de la base elegida.

\textbf{Este libro no extrae de ahí ninguna afirmación sobre la población indígena del
departamento}, que requeriría un procesamiento distinto y que el capítulo
\ref{cap:resistencias} aborda por otra vía.
\end{cautela}

\section{Qué se constituye hoy en el departamento}

Todo lo anterior mira hacia atrás: un censo de 2018, unos decretos de 2014 y 2015. Conviene
cerrar mirando lo que se está constituyendo ahora, porque las sociedades publican su acto
constitutivo y la publicación dice el objeto, el capital y el domicilio.

Entre julio de 2025 y marzo de 2026 se publicó la constitución de tres sociedades con sede en el
departamento.

\begin{ficha}{Boletín Oficial Nº 22.147, 12 de marzo de 2026 --- \textsc{agro san alejo s.a.s.}}
Constituida el 15 de enero de 2026. \textbf{Sede social: Finca San Alejo, matrícula Nº 1.539},
Ruta Nacional Nº 9, altura Dique Campo Alegre, departamento La Caldera.

\textbf{Objeto}, en el orden en que lo enumera el instrumento: \textbf{1.\ Agroturismo y turismo
rural} ---alojamientos rurales, casas de campo, posadas, campings, servicios gastronómicos,
eventos y ferias---; \textbf{2.\ Producción agropecuaria y forestal}; 3.\ Servicios ambientales
y sustentabilidad; 4.\ Comercialización; \textbf{5.\ Actividades inmobiliarias}, incluida «la
realización de desarrollos inmobiliarios vinculados al objeto social».

Capital: \$810.000. \textbf{Los dos socios tienen domicilio fuera de la provincia}: Ciudad
Autónoma de Buenos Aires y partido de Pilar.
\end{ficha}

\textbf{San Alejo es una de las fincas que el artículo 6º del PDUA enumera}, según el capítulo
\ref{cap:pdua}, y aparece en la cadena que el capítulo \ref{cap:fincas} sigue desde 1909. Su
matrícula ---la 1.539--- es nueva para este libro.

\begin{ficha}{Boletín Oficial Nº 22.148, 13 de marzo de 2026 --- \textsc{cantera el rescoldo s.a.}}
Constituida el 27 de enero de 2026. Sede social en la ruta 9, kilómetro 1627, La Caldera.

\textbf{Objeto}: exploración y explotación de minas y canteras; «\textit{depósitos de canto
rodado, arenas, arcillas y similares}»; \textbf{energía} ---térmica, hidroeléctrica,
fotovoltaica, eólica y renovable---; mandatos, comercialización, industrialización e
importación.

\textbf{Capital: \$30.000.000.} Dos socios, con domicilio en San Lorenzo y en la ciudad de
Salta; uno de ellos declara la profesión de \textbf{martillero público y corredor
inmobiliario}.
\end{ficha}

\begin{ficha}{Boletín Oficial Nº 21.994, 22 de julio de 2025 --- \textsc{campo alegre constructora s.r.l.}}
Constituida el 30 de mayo de 2025. Sede social sobre la Ruta Nacional 9, localidad de La
Caldera.

\textbf{Objeto}: construcción de viviendas, obras civiles y \textbf{desarrollos inmobiliarios},
con mención expresa de dúplex, departamentos y sistemas industrializados en seco.

Capital: \$1.500.000. Tres socios; dos declaran domicilio en el departamento y uno en Buenos
Aires.
\end{ficha}

\begin{cautela}
\textbf{Qué prueban estos tres instrumentos y qué no.}

\textbf{Prueban} que las tres sociedades existen, con esa sede, ese objeto y ese capital. No
prueban que operen, que hayan empezado ninguna actividad, ni ---en el caso de la cantera---
que tenga concesión: el capítulo \ref{cap:resistencias} muestra que las concesiones de áridos
las otorga el Juzgado de Minas por un trámite distinto, y en la nómina oficial de 2022 esta
sociedad no podía figurar porque no existía.

\textbf{Y no dicen nada sobre las personas.} Este libro consigna los domicilios porque son un
dato estructural ---de dónde viene el capital que se radica en el departamento--- y no
construye nada sobre los socios, que constituyen sociedades como cualquiera puede hacerlo.

\textbf{Una inconsistencia que conviene registrar:} los tres instrumentos ubican la sede sobre la ruta 9, pero no coinciden en cómo la denominan, y sólo uno consigna el kilometraje. Las fichas de este apartado normalizan la denominación de la ruta.
\end{cautela}

\textbf{Lo que sí se puede leer en los tres, puestos juntos, es la composición.} Una sociedad
sobre una finca histórica cuyo objeto pone \textbf{el agroturismo antes que la producción} y
termina en desarrollos inmobiliarios. Una constructora de viviendas y dúplex. Y una sociedad de
canteras y energía con un capital treinta y siete veces mayor que el de la primera.

\textbf{Es exactamente la composición que este capítulo describió para el período
2002--2018, constituyéndose otra vez en 2025 y 2026.} La agropecuaria aparece, pero en segundo
lugar y dentro de un objeto que empieza por el turismo.

\section{¿Fue por diseño?}

Todo lo anterior describe una transformación sin nombrar a nadie que la haya decidido. Queda
por examinar la posibilidad contraria: que en algún documento conste que este departamento
debía dejar de ser agropecuario para volverse turístico.

Si ese documento existe, hay uno solo donde podría estar.

El capítulo \ref{cap:pdua} ya lo examinó y transcribió el pasaje decisivo: el plan declara que
el crecimiento provincial fluyó «espontáneamente, \textbf{sin planificación}», nombra como
problema visible el «\textbf{dinámico crecimiento del área metropolitana}» y dice que hay que
«\textbf{paliar los efectos de la falta de previsión}».

\textbf{El instrumento de planificación que la provincia aprobó con carácter definitorio, invocando el artículo 77 de su Constitución, declara, en
2012, que el crecimiento del área metropolitana ocurrió sin planificación.} No asigna a La
Caldera un perfil turístico ni ningún otro: registra que lo que pasó ahí pasó solo.

Y conviene agregar dos datos de estructura que ese capítulo no necesitaba.

\textbf{El plan se organiza en trece sectores} ---salud, educación, agropecuario, minería,
industria, comercio, energía, turismo, ambiente, planeamiento urbano, infraestructura,
tecnología y economía--- y siete ejes provinciales. \textbf{No tiene estructura territorial por
departamento.}

\textbf{Y la Intendencia de La Caldera participó} de la etapa de Diagnóstico: figura en la
nómina de instituciones convocadas que asistieron.

\subsection*{La asimetría}

Eso podría leerse como un límite propio del género: los planes estratégicos son generales y no
bajan a departamentos. \textbf{Pero este sí baja, y en el mismo documento.}

El capítulo de Minería del PDES 2030 \textbf{cartografía los salares de la puna salteña}
---departamento Los Andes---, cuantifica sus recursos de litio, boro y sulfato de sodio, y
\textbf{nombra a las empresas que los exploraban}, con su origen y sus cifras.

\begin{cautela}
\textbf{Ésa es la observación de este apartado, y es toda.}

El mismo plan, con el mismo rango, \textbf{territorializa cuando hay litio y no territorializa
el corredor metropolitano}. Para la puna nombra el recurso, el inversor y la condición; para el
área metropolitana escribe que el crecimiento fluyó espontáneamente y que hay que paliar los
efectos de la falta de previsión.

\textbf{De modo que la respuesta a la pregunta del título es negativa, y de una manera
particular.} La Caldera no fue diseñada como departamento turístico. \textbf{Quedó fuera del
diseño.} Lo que este capítulo documenta ---la desaparición de la superficie agropecuaria, la
extracción para afuera, la vivienda de fin de semana, los contratos de promoción con exención---
ocurrió en un territorio que el único plan definitorio de la provincia describió como un lugar
donde las cosas venían pasando sin que nadie las previera.

\textbf{Y eso no exime: agrava.} Porque el plan lo dijo en 2012, y el capítulo
\ref{cap:opacidad} muestra que el criterio siguió sin escribirse durante los catorce años
siguientes.
\end{cautela}

\begin{cautela}
\textbf{Dos advertencias sobre este apartado, que conviene que sean explícitas.}

\textbf{La primera es de alcance, y quedó saldada.} Este relevamiento consultó primero las
ciento cincuenta páginas iniciales del Informe Final y declaró que faltaban los capítulos de
Planeamiento Urbano Territorial, Medio Ambiente y Turismo. \textbf{Después los leyó.} El
documento completo ---ciento setenta páginas--- se recorrió entero.

\textbf{«La Caldera» aparece dos veces en todo el Plan de Desarrollo Estratégico de la Provincia
de Salta 2030.} Una, en la lista de instituciones que participaron del Diagnóstico:
«\textit{Intendencia de La Caldera}». Y otra, en un inventario de establecimientos hoteleros del
capítulo de Turismo, donde figura con \textbf{tres hoteles y parahoteles} ---ninguno
categorizado como hotel; el inventario nombra una hostería y una casa de campo---, contra seis de Vaqueros,
veinticinco de San Lorenzo y ciento ochenta y siete de la ciudad de Salta.

\textbf{Eso es todo.} No hay asignación de perfil, ni zonificación, ni mención del corredor. La
expresión «área metropolitana» aparece dos veces en el documento entero, y una de ellas es el
pasaje que el capítulo \ref{cap:pdua} transcribe: el que dice que creció sin planificación.

\textbf{La segunda es de interés.} La consultora del sector Educación del PDES 2030 fue la
\textbf{Prof.\ Constanza Diedrich}, \textbf{hermana del autor de este libro}. El sector
educativo no interviene en ninguna parte de la cadena argumentativa de este capítulo ni de este
libro, y la mención del plan se limita a su estructura, a su carácter normativo y a dos
párrafos de su área territorial. Se declara igualmente, porque un trabajo que exige a los
organismos que hagan explícito su criterio no puede reservarse los propios.
\end{cautela}

\section{Qué afirma este capítulo y qué no}

\begin{cautela}
\textbf{Afirma} que la superficie agropecuaria del departamento cayó un 88,4 por ciento entre
2002 y 2018 mientras la provincial crecía; que había producción reconocida por el Estado ---
tabaco en 1980, olivo promovido en 1934, un aserradero ---; que las dos actividades que
crecieron en el período extraen o alojan para afuera; y que el Estado provincial ratificó por
decreto dos contratos de promoción turística con exenciones cuyo contenido no publicó.

\textbf{No afirma} que exista relación causal entre esas cosas. No hay un acto que desaliente
la producción agropecuaria y aliente la turística: hay dos series que corren en paralelo, y
este libro las pone una al lado de la otra porque nadie lo había hecho.

\textbf{Y no afirma que el turismo sea el problema.} Un departamento a veinte kilómetros de una
capital de medio millón de habitantes, con dique, cerros y clima, iba a tener turismo con
promoción o sin ella. La pregunta que queda abierta no es si debió haberlo, sino
\textbf{por qué el Estado documentó con tanto detalle lo que promovía y con tan poco lo que
dejaba de promover}.
\end{cautela}

